\section{Who can see a beat}
\label{sec:whosees}

Return to the click trains. Two trains of clicks in the air \emph{add}: the
pressure at your eardrum is the sum of the two. Two gratings printed on
transparencies and held to the light \emph{multiply}: light gets through
where both are clear, and the transmission is the product. Two gratings
thrown on a wall by two projectors add again. The difference sounds like
bookkeeping and is the whole reason a beat in the air needs an ear and a
beat on paper does not.

On the torus the two cases are two pictures. If family 1 has profile $g_1$
(one where a click or a stroke is, zero between) and family 2 has $g_2$,
then an additive superposition has the picture $I(x_1,x_2)=g_1(x_1)+g_2(x_2)$
and a multiplicative one has $I(x_1,x_2)=g_1(x_1)\,g_2(x_2)$. Everything
follows from expanding each in waves.

\begin{lookup}{the waves of a pulse train}
What to look up: \emph{the Fourier coefficients of a rectangular pulse
train} (or ``Fourier series of a square wave with duty cycle $d$''). The
consequences used: a train that is on for a fraction $d$ of each period has
coefficients $\hat g(n)=\sin(n\pi d)/(n\pi)$ for $n\ne0$ and $\hat g(0)=d$;
so $\hat g(n)$ is exactly zero when $nd$ is a whole number (a $50\%$ train
has no even harmonics at all), and otherwise $|\hat g(n)|$ falls like
$1/n$.
\end{lookup}

\begin{theorem}[Who sees a beat]
\label{thm:whosees}
Write a recipe $k=(a,b)$ as \emph{cross} when both $a$ and $b$ are nonzero.
\begin{enumerate}
\item An additive picture $g_1(x_1)+g_2(x_2)$ has \emph{no cross recipes}:
its waves are $(n,0)$ and $(0,n)$ only. A linear pooling observer, one that
averages the superposition and nothing more, sees no beat in an additive
superposition, ever.
\item A multiplicative picture $g_1(x_1)\,g_2(x_2)$ carries every cross
recipe with strength $\hat I(a,b)=\hat g_1(a)\,\hat g_2(b)$: a linear
observer sees the beat, with the strength of the two harmonics that make it.
\item A front end that bends, $N''\ne0$, mints cross recipes out of an
additive picture. Squaring the sum gives
$(g_1+g_2)^2=g_1^2+2g_1g_2+g_2^2$, whose middle term is a product: the
difference recipe $(1,-1)$ appears with strength $\hat g_1(1)\hat g_2(1)/2$
when the sum is normalised to $(g_1+g_2)/2$.
\end{enumerate}
\end{theorem}

\begin{proof}
A wave on the torus with recipe $(a,b)$ is $e^{2\pi i(ax_1+bx_2)}$. The
expansion of $g_1(x_1)$ contains only waves $e^{2\pi i\,n x_1}$, recipes
$(n,0)$, and likewise $g_2$; a sum of the two contains nothing else, which
is the first claim, and by Theorem~\ref{thm:sweep} any average that removes
the single-family recipes removes everything. The product of two expansions
is the sum over all pairs of the products of waves, and
$e^{2\pi i\,ax_1}e^{2\pi i\,bx_2}$ is the wave of recipe $(a,b)$ with
coefficient $\hat g_1(a)\hat g_2(b)$, which is the second. Expanding the
square gives the third, the cross term $2g_1g_2$ contributing
$2\hat g_1(a)\hat g_2(b)$ at $(a,b)$, and the normalisation divides by four.
\end{proof}

Measured: the additive beat under a linear observer is $\obsAdditiveBeat$,
zero to the precision of the arithmetic; the printed beat is
$\obsPrintedBeat$; a squaring front end mints the additive beat at
$\obsSquaringBeat$, exactly half the printed value; a saturating front end
$\min(g_1+g_2,1)$, the front end of a photoreceptor or a clipping amplifier,
mints it at $\obsSaturatingBeat$, nearly the printed value
(script \texttt{observer.mjs}).

This is the coincidence detector of \S\ref{sec:coincidence}, explained. A
detector that fires only when both clicks are present is a front end that
multiplies --- ``both'' is a product --- so its picture on the torus is the
small square where both counts are near a whole number, a picture with
every cross recipe in it, and the slowest cross recipe is the difference of
counts $D$. The detector's output is a wave in $D$, once per unit of $D$,
once per second, and nothing else survives its pooling. It is also
Helmholtz's combination tone~\cite{Helmholtz1954Sensations}: two tones in
the air have no difference tone between them until an ear, which is not
linear before it pools, makes one. The distinction is not between kinds of
observer. It is between kinds of superposition, and the theorem puts it in
the only place it can be: which recipes the observer's picture contains.

\subsection{Hard patterns are observer-proof, soft ones are not}
\label{sec:hardsoft}

\begin{corollary}[Hard patterns]
\label{cor:hard}
If the picture $I$ takes only two values, then for any front end $N$,
\[
  N\circ I \;=\; N(0)+\big(N(1)-N(0)\big)\,I ,
\]
a constant plus a multiple of $I$: every observer sees the same recipes at
the same relative strengths, and differs from every other only by a
constant and a gain. A hard-edged pattern --- ink or no ink, click or no
click --- has one third pattern, and every observer sees it.
\end{corollary}

The proof is the display: a function of a variable that takes two values is
determined by its two values, and any such function is affine. Measured: the
octave null of \S\ref{sec:selection} under a linear and a squaring observer
agree to $\obsHardNull$.

\begin{corollary}[Soft patterns]
\label{cor:soft}
Blur the edges of a $50\%$ train symmetrically, by any even smoothing. Its
even harmonics stay exactly zero, so every \emph{linear} observer keeps the
null at every blur. A squaring observer does not: it reopens the null with
a strength that grows linearly in the width of the blur.
\end{corollary}

\begin{proof}
A $50\%$ train satisfies $g(x+\tfrac12)=1-g(x)$; a symmetric smoothing
preserves this half-turn antisymmetry, and a function with it has no even
harmonics. Squaring destroys the antisymmetry exactly on the blurred edges
and nowhere else, so the even harmonics it mints scale with the edge width.
\end{proof}

Measured on strokes with anti-alias ramps of half-width \obsReopenRamps\
units: the linear null holds at every ramp, and squaring reopens it at
$\obsReopenAmps$, linear in the ramp (\texttt{observer.mjs}).

\subsection{The same theorem in an ear}
\label{sec:ear}

Everything so far was measured in a drawing tool. The theorem does not know
that. Take two pulse trains at $200$ and $405$ clicks per second --- a
mistuned octave, whose slow recipe is $(2,-1)$, twice the lower count minus
the upper, beating five times a second --- add them as sounds add, and
listen through four model ears: a low-pass filter alone (a linear ear),
a squaring front end followed by the low-pass (the textbook square-law
detector), a cubing front end, and a two-stage ear that squares, pools,
and squares again. Sweep the lower train's duty from $30\%$ to $70\%$ and
read the five-hertz line at the output. Figure~\ref{fig:ear} is the result
(script \texttt{ear.mjs}).

\begin{figure}[t]
  \centering
  \begin{tikzpicture}
    \begin{axis}[width=0.5\linewidth, height=5.4cm, ymode=log,
      title={hard pulses}, xlabel={duty of the lower train},
      ylabel={the 5 Hz line at the ear's output},
      xmin=0.27, xmax=0.73, ymin=1e-7, ymax=1,
      xtick={0.3,0.4,0.5,0.6,0.7},
      tick align=outside, tick pos=left, axis line style={gray!60},
      grid=major, grid style={gray!18, very thin},
      label style={font=\footnotesize}, tick label style={font=\scriptsize},
      title style={font=\footnotesize},
      legend style={font=\scriptsize, draw=gray!40, at={(0.5,-0.32)}, anchor=north, legend columns=2}]
      \addplot[soft, mark=x, line width=0.6pt] table[x=duty, y=hard_linear, col sep=comma]{data/ear-octave.csv};
      \addlegendentry{linear}
      \addplot[accent, mark=*, mark size=1.2pt, line width=0.7pt] table[x=duty, y=hard_squarelaw, col sep=comma]{data/ear-octave.csv};
      \addlegendentry{square-law}
      \addplot[cool, mark=square*, mark size=1.1pt, line width=0.7pt] table[x=duty, y=hard_cubic, col sep=comma]{data/ear-octave.csv};
      \addlegendentry{cubic}
      \addplot[warm, mark=triangle*, mark size=1.3pt, line width=0.7pt] table[x=duty, y=hard_squarepoolsquare, col sep=comma]{data/ear-octave.csv};
      \addlegendentry{square, pool, square}
    \end{axis}
  \end{tikzpicture}%
  \hfill%
  \begin{tikzpicture}
    \begin{axis}[width=0.5\linewidth, height=5.4cm, ymode=log,
      title={softened pulses}, xlabel={duty of the lower train},
      xmin=0.27, xmax=0.73, ymin=1e-7, ymax=1,
      xtick={0.3,0.4,0.5,0.6,0.7}, yticklabels={},
      tick align=outside, tick pos=left, axis line style={gray!60},
      grid=major, grid style={gray!18, very thin},
      label style={font=\footnotesize}, tick label style={font=\scriptsize},
      title style={font=\footnotesize}]
      \addplot[accent, mark=*, mark size=1.2pt, line width=0.7pt] table[x=duty, y=soft_squarelaw, col sep=comma]{data/ear-octave.csv};
      \addplot[cool, mark=square*, mark size=1.1pt, line width=0.7pt] table[x=duty, y=soft_cubic, col sep=comma]{data/ear-octave.csv};
      \addplot[warm, mark=triangle*, mark size=1.3pt, line width=0.7pt] table[x=duty, y=soft_squarepoolsquare, col sep=comma]{data/ear-octave.csv};
    \end{axis}
  \end{tikzpicture}
  \caption{\figlead{The octave beat through four ears.} Two pulse trains at
  $200$ and $405$ per second, added as sounds add, through a linear ear, a
  square-law ear, a cubic ear and a two-stage ear, as the lower train's duty
  is swept. Left, hard-edged pulses: the linear ear hears nothing at any
  duty (Theorem~\ref{thm:whosees}, first claim), and every nonlinear ear
  hears the beat except at duty one half, where it is extinguished
  $\earHardNullSquare\times$, $\earHardNullCubic\times$ and
  $\earHardNullCascade\times$ under the three ears --- a hard pattern is
  observer-proof (Corollary~\ref{cor:hard}). Right, the same pulses with
  softened edges: the square-law ear keeps the null ($\earSoftNullSquare\times$),
  the cubic ear reopens it to within $\earSoftNullCubic\times$ of its
  neighbours (Corollary~\ref{cor:soft}, where cubing plays the part of
  squaring one order up). Script \texttt{ear.mjs}, data \texttt{ear.json}.}
  \label{fig:ear}
\end{figure}

Three things in the figure were predictions before they were curves. The
linear ear hears no beat at any duty, $\earLinearOverSquare$ of what the
square-law ear hears, because a sum of trains has no cross recipe. Every
nonlinear ear loses the beat at duty one half, because the beat rides the
lower train's second harmonic and a $50\%$ train has none: the same duty
null that \S\ref{sec:selection} measures on paper, in sound, through three
different ears at once. And softening the pulses separates the ears: a
square-law front end's cross term is a product of one harmonic from each
train, so it inherits the lower train's missing second harmonic and keeps
the null; a cubic front end's $g_1^2g_2$ term carries the \emph{square} of
the lower train, whose second harmonic is not zero once the edges are soft,
and the null reopens. That is a prediction about people: \emph{which
nonlinearity an ear has is audible}, in a mistuned octave of blurred square
waves, as the presence or absence of a five-hertz beat. We have not done the
listening test. \S\ref{sec:ledger} says exactly what it would take.

\begin{tryit}{a beat that a duty can switch off}
Any synthesizer with a pulse-width control, or a free web synth. Two
oscillators, pulse waves, at $200$ and $405$ hertz; set the lower one's
width to $30\%$ and you hear a five-per-second throb; set it to $50\%$ and
the throb goes, though both tones are still there; $70\%$ and it is back.
Your ear is the nonlinearity. Then put a gentle low-pass filter on each
oscillator, so the edges soften, and listen at $50\%$ again: if you hear the
throb come back, your ear's front end is not a square-law.
\end{tryit}
